\section{Discussion and Future Work}
\label{sec:discussion}

\paragraph{What EWAT is, and is not.}
EWAT is a working early-warning pipeline whose honest, independently-validated contribution is
active-scenario typing at $0.920$ macro-AUROC with confirmed temporal precursion
($\Delta=-0.116$). It is not a generic event predictor: the learned encoder adds geometric and
ontological value, not aggregate predictive power, and open-set generalisation to unseen fault
types remains only partially solved. We consider the negative and nuanced results --- a robustly
falsified look-through, a circular metric demoted in favour of an independent one, an unstable type
count --- to be as much a contribution as the positive headline, because they delimit precisely
where the approach is and is not trustworthy.

\paragraph{Roadmap.}
Four research axes follow directly from the limitations.
\emph{Axis~A (ontology$\,\leftrightarrow\,$prediction coupling)} would feed the causal graph back
into the precursor as upstream features or a post-hoc re-ranker, testing whether the ontology adds
predictive value rather than only descriptive value.
\emph{Axis~B (robust temporal precursion)} would replace the STGCN with a temporal transformer,
multi-seed the C2 distant-window result, and validate lead time cross-dataset on RCAEval.
\emph{Axis~C (open-set)} would extend the partial OpenMax result with Mahalanobis-based and
energy-based detectors and few-shot meta-learning.
\emph{Axis~D (deployment)} would specify the streaming pipeline, the retraining trigger, and a
formal alert contract.

\paragraph{Towards \texttt{ewat\_v5}.}
The most consequential next step is data. \texttt{ewat\_v5} abandons the six-service Online Boutique
for \emph{Train~Ticket} (41 Spring~Cloud services), with a much larger episode budget and, crucially,
a held-out set drawn from the project's documented \emph{real} bugs (the F1--F12 catalogue) ---
a literature gold standard for RCA-style evaluation. This directly attacks L1 (episode length),
L13 (topology scale), L14 (cross-system validity), and L17 (real rather than synthetic faults),
and would let us re-test every hypothesis on a substantially harder and more realistic system.
